\chapter*{Presentación}
\markboth{Presentación}{}
\addcontentsline{toc}{chapter}{Presentación}
Este volumen desarrolla el primer curso de una colección de cuatro libros de análisis. Su hilo conductor es el paso de la distancia a la estructura lineal y, finalmente, a la geometría del producto interno. La intención es que un estudiante de la mitad de la licenciatura en Matemáticas pueda seguir las pruebas y entender por qué se eligen sus pasos, sin rebajar la precisión de los enunciados.

El punto de partida son las 32 lecturas de \emph{Análisis Matemático 1, Licenciatura en Física y Matemáticas}, de Miguel Ángel Valencia Bucio \cite{valencia}. Las lecturas se integran en seis partes y dieciséis capítulos temáticos. Se conserva la identificación de los ejercicios originales y una tabla final permite localizar el contenido de cada lectura. La autoría de esas lecturas se reconoce expresamente: la organización editorial, las explicaciones ampliadas, las soluciones desarrolladas y las rectificaciones forman esta edición de Eder Said Martinez Ramirez.

Un resultado falso no puede completarse añadiendo pasos a una prueba. Cuando un enunciado necesita una hipótesis o admite un contraejemplo, el texto explica el problema y formula una versión válida. Se distinguen asimismo las erratas matemáticas de las pérdidas de barras, índices o signos ocasionadas por la extracción del PDF. Las soluciones incluyen el enunciado que se utiliza y señalan las correcciones pertinentes.

\section*{Créditos de las ampliaciones y del diseño}
Los apuntes de Egor Maximenko, profesor de la ESFM del IPN, orientan las ampliaciones sobre espacios cocientes y seminormas, el teorema de Picard, matrices de Gram y proyecciones, y bases de Schauder \cite{max-cociente,max-seminormado,max-picard,max-gram,max-proyeccion,max-schauder}. Cada bloque identifica sus documentos concretos y cada ejercicio adaptado señala su procedencia. La redacción ampliada, las comprobaciones, las soluciones y los ejercicios adicionales elaborados para esta edición se distinguen de los enunciados tomados como punto de partida. El volumen conserva los 206 ejercicios de las lecturas e incorpora 20 ejercicios de ampliación, todos acompañados de solución.

El diseño toma como referencia visual los \emph{Apuntes de Homotopía Estable} del repositorio público de MrPaborru (\url{https://github.com/MrPaborru/Apuntes-Homotop-a-Estable}), combinados con una composición tradicional de libros de matemáticas. El emblema de la ESFM procede de su sitio oficial. Su uso identifica la adscripción académica del autor; esta colección es una edición personal, no una publicación oficial ni una obra avalada por el IPN. Se conserva una paleta inspirada en los colores de la ESFM y una tipografía Latin Modern. Los detalles de procedencia gráfica se documentan también en las fuentes LaTeX.

\section*{Cómo estudiar este volumen}
Conviene conocer operaciones con conjuntos, cuantificadores, inducción, sucesiones y series reales, cálculo de una variable y álgebra lineal elemental. No se presupone un curso previo de topología ni de análisis funcional. La completitud de los números reales y los resultados básicos del cálculo se usan como conocimientos previos. El lema de Zorn se enuncia cuando interviene como principio de existencia; las construcciones que dependen de él se separan de las construcciones explícitas.

Antes de leer una demostración, identifique sus hipótesis y escriba lo que se quiere producir: un radio, una subsucesión, una cota o un operador. Lea después la estrategia y trate de reconocer dónde se usa cada hipótesis. En los ejercicios, conviene intentar una solución antes de consultar el desarrollo. Un contraejemplo merece el mismo cuidado que una prueba: hay que verificar tanto las hipótesis que satisface como la conclusión que contradice.

La jerarquía de la colección es volumen, parte, capítulo y sección. Cada volumen corresponde a un curso; sus partes reúnen etapas conceptuales y sus capítulos desarrollan unidades matemáticas. En este volumen, las seis partes recorren espacios métricos y continuidad, completitud, compacidad y estructura topológica, espacios normados y operadores, espacios de funciones y geometría de Hilbert. Los capítulos se numeran consecutivamente dentro del volumen.

Los ejercicios se distribuyen cerca de las herramientas que utilizan. Su número dentro del volumen permite citarlos sin ambigüedad; la indicación entre paréntesis conserva la lectura y el número de origen. Las referencias enlazadas llevan al resultado o ejercicio correspondiente. La tabla final de correspondencia distingue el orden temático del libro del orden de las lecturas originales.

\section*{Convenciones y fuentes}
Salvo indicación contraria, $\N=\{1,2,\ldots\}$ y $\F$ es $\R$ o $\C$. El producto interno es lineal en su primer argumento. Las bolas abiertas y cerradas se distinguen del cierre de una bola abierta: en un espacio métrico arbitrario esos dos últimos conjuntos pueden ser diferentes. El dual $V^*$ contiene solamente los funcionales lineales continuos.

Rudin \cite{rudin1976} ofrece un desarrollo de topología métrica, sucesiones de funciones y aproximación; Munkres \cite{munkres2000} amplía el contexto topológico; Kreyszig \cite{kreyszig1978} y Conway \cite{conway1990} permiten continuar hacia espacios normados, Hilbert y operadores. La bibliografía final identifica obras reales y ediciones comprobadas. Las referencias externas son orientación y contraste; las demostraciones desarrolladas aquí deben poder seguirse a partir de los resultados previos indicados, sin buscar una página no especificada en otro libro.
\section*{Apuntes y solucionario}
Los ejercicios se presentan en estos apuntes sin soluciones. Los mismos
enunciados, con la misma numeración por capítulo y sus soluciones, se reúnen
en el documento separado \emph{Solucionario} de este volumen. Las menciones
a las soluciones y a los ejercicios resueltos de esta presentación se refieren
a ese documento complementario.

